% !TEX program = xelatex
\documentclass[11pt,a4paper]{article}

% XeLaTeX packages for Unicode support
\usepackage{fontspec}
\usepackage{xeCJK}
\setCJKmainfont{HaranoAjiMincho}[Scale=0.9]
\setmainfont{Latin Modern Roman}

\usepackage[margin=1in]{geometry}
\usepackage{hyperref}
\usepackage{titlesec}
\usepackage{enumitem}
\usepackage{fancyhdr}
\usepackage{lastpage}
\usepackage{datetime2}
\usepackage{booktabs}
\usepackage{array}

% Colors
\usepackage{xcolor}
\definecolor{linkcolor}{RGB}{0,51,102}
\definecolor{darknavy}{RGB}{0,32,76}
\definecolor{lightgray}{RGB}{180,180,180}

% Hyperref settings
\hypersetup{
    colorlinks=true,
    linkcolor=linkcolor,
    urlcolor=linkcolor,
    pdfauthor={Yuta Takahashi},
    pdftitle={研究業績 - 高橋悠太}
}

% Section formatting
\titleformat{\section}
    {\large\bfseries\scshape\color{darknavy}}
    {}
    {0em}
    {}[\titlerule]
\titlespacing{\section}{0pt}{16pt}{8pt}

% Header/Footer
\pagestyle{fancy}
\fancyhf{}
\renewcommand{\headrulewidth}{0pt}
\fancyfoot[C]{{\small\color{gray} \thepage\ / \pageref{LastPage} ページ}}

% Custom commands
\newcommand{\entrytitle}[1]{{\bfseries\color{darknavy}#1}}

\begin{document}

% Header
\begin{center}
    {\LARGE\bfseries\color{darknavy} 高橋 悠太}\\[2pt]
    {\large (たかはし ゆうた / Yuta Takahashi)}\\[8pt]
    {\color{lightgray}\rule{0.6\textwidth}{0.4pt}}\\[8pt]
    〒567-0047　大阪府茨木市美穂ヶ丘6-1\\[2pt]
    大阪大学 社会経済研究所\\[4pt]
    \href{mailto:yuta.takahashi@iser.osaka-u.ac.jp}{yuta.takahashi@iser.osaka-u.ac.jp}\\[10pt]
    {\small\color{gray} 作成日：\today}
\end{center}

\vspace{8pt}

%----------------------------------------------------------------------
\section{（１）氏名}

\noindent\textbf{高橋 悠太}（たかはし ゆうた / Yuta Takahashi）

\vspace{8pt}

%----------------------------------------------------------------------
\section{（２）研究の概要}

\noindent
マクロ経済学を専門とし、日本経済が直面する主要課題（財政赤字・低成長・少子高齢化）を中心に研究を行っている。
具体的には、インフレ期待の形成メカニズム、ゼロ金利制約下での金融政策、国際貿易における重力モデル、技術進歩の停滞、
名目賃金の下方硬直性などを研究対象としている。理論・実証の両面からアプローチし、
日本および世界経済の停滞・スタグフレーション現象の解明にも取り組んでいる。

\vspace{8pt}

%----------------------------------------------------------------------
\section{（３）学歴・職歴}

\subsection*{学歴}
\vspace{4pt}
\begin{tabular}{@{}lll@{}}
    2009年 & 京都大学 経済学部 & 学士（経済学）\\[4pt]
    2011年 & 京都大学大学院 経済学研究科 & 修士（経済学）\\[4pt]
    2018年 & Northwestern University & Ph.D. in Economics\\
\end{tabular}

\vspace{12pt}

\subsection*{職歴}
\vspace{4pt}
\begin{tabular}{@{}lll@{}}
    2018年9月 -- 2020年3月 & 一橋大学 経済研究所 & 非常勤研究員（ポスドク）\\[4pt]
    2020年4月 -- 2025年8月 & 一橋大学 経済研究所 & 准教授\\[4pt]
    2025年9月 -- 現在       & 大阪大学 社会経済研究所 & 准教授（任期なし）\\
\end{tabular}

\vspace{8pt}

%----------------------------------------------------------------------
\section{（４）主な発表論文名・著書名等}

\subsection*{査読付き論文}
\vspace{4pt}
\begin{enumerate}[leftmargin=*, labelsep=8pt, itemsep=8pt]

    \item \textbf{``Universal Gravity''}\\
    {\small Treb Allen, Costas Arkolakis との共著}\\
    \textit{Journal of Political Economy}, 128(2): 393--433, 2020

    \item \textbf{``A Parsimonious Model for Zero Inflation at the Zero Lower Bound''}\\
    {\small Naoki Takayama との共著}\\
    \textit{Economics Letters}, Volume 247, 2025年2月

    \item \textbf{``An Experiment on a Multi-Period Beauty Contest Game''} (2025)\\
    {\small Nobuyuki Hanaki との共著}\\
    \textit{Forthcoming at Experimental Economics}

\end{enumerate}

\subsection*{改訂中論文}
\vspace{4pt}
\begin{enumerate}[leftmargin=*, labelsep=8pt, itemsep=8pt]

    \item \textbf{``Robustly Optimal Voting Rule''} (2025)\\
    {\small Noriaki Kiguchi, Shinpei Noguchi との共著}\\
    \textit{Conditionally accepted at Journal of Economic Theory}

\end{enumerate}

\subsection*{ワーキングペーパー}
\vspace{4pt}
\begin{enumerate}[leftmargin=*, labelsep=8pt, itemsep=8pt]

    \item \textbf{``Does Expected Inflation Matter? Evidence from Japanese Value-Added Tax Hikes''} (2024)\\
    {\small Naoki Takayama との共著}

    \item \textbf{``Global Technology Stagnation''} (2024)\\
    {\small Naoki Takayama との共著}

    \item \textbf{``Hidden Stagflation''} (2024)\\
    {\small Naoki Takayama との共著}

    \item \textbf{``Anchoring Inflation Expectation''} (2020)\\
    {\small Lawrence Christiano との共著}

\end{enumerate}

\vspace{8pt}

%----------------------------------------------------------------------
\section{（５）その他}

\subsection*{学会・セミナー・講演}
\vspace{4pt}
\begin{itemize}[leftmargin=*, itemsep=4pt]
    \item \textbf{Seminar in Macroeconomics} 共同オーガナイザー\\
    一橋大学 経済研究所・京都大学 経済研究所の共同運営セミナー
\end{itemize}

\vspace{8pt}

\subsection*{受賞歴}
\vspace{4pt}
\begin{itemize}[leftmargin=*, itemsep=4pt]
    \item （追記予定）
\end{itemize}

\vspace{8pt}

\subsection*{学術誌レフェリー等}
\vspace{4pt}
\begin{itemize}[leftmargin=*, itemsep=4pt]
    \item （追記予定）
\end{itemize}

\end{document}
